\documentclass[11pt]{article}
\usepackage{amsmath, amssymb, amsthm}
\usepackage{hyperref}
\usepackage{geometry}
\geometry{margin=1in}

\title{Fun-Sphere: A Holonomy-Induced Submanifold for Reversible Narrative Ecologies}
\author{Borealis S. Hedling}
\date{October 2026}

\begin{document}
\maketitle

\begin{abstract}
This preprint introduces the Fun-Sphere, a holonomy-induced submanifold
embedded inside a neutral high-dimensional geometric substrate. The
Fun-Sphere arises when a reversible operator loop returns a point to
itself while producing a non-identity frame transform. This phenomenon
corresponds to non-trivial holonomy and defines a sovereign, reversible,
structurally isolated region of the manifold.

We demonstrate how this geometric construction provides a toy-model
foundation for reversible narrative ecologies, where story states evolve
through operator algebra rather than branching logic. The Fun-Sphere
serves as a mathematically grounded ``pocket universe'' suitable for
simulation-based storytelling, reversible puzzles, and emergent
narrative physics.
\end{abstract}

\section{Introduction}

Traditional narrative engines rely on branching structures or scripting
languages. These approaches scale poorly and fail to capture reversible
or emergent story behavior. This work explores a geometric alternative:
treat narrative as manifold traversal.

Operators correspond to story actions, parallel transport corresponds to
meaning propagation, holonomy corresponds to narrative curvature, and
submanifolds correspond to sovereign story worlds. The Fun-Sphere is the
simplest non-trivial example of such a world.

\section{Ambient Manifold}

Let $M$ be an $n$-dimensional neutral manifold equipped with the identity
metric:
\[
g = I_n.
\]
Points in $M$ are represented as vectors $x \in \mathbb{R}^n$. The
manifold is intentionally structure-free: no curvature, no torsion, and
no coordinate dependencies.

\section{Reversible Operators}

A reversible operator is defined as a pair $(T, T^{-1})$ where
$T \in GL(n)$. Operators act on points via:
\[
x' = T x.
\]
A sequence of operators $\gamma = \{T_1, \dots, T_k\}$ composes into:
\[
T_{\gamma} = T_k \cdots T_1.
\]

\section{Parallel Transport}

Parallel transport is purely algebraic:
\[
P_{\gamma}(x) = T_{\gamma} x.
\]
This defines a flat connection with no Christoffel symbols or curvature
tensors.

\section{Holonomy}

Holonomy is defined as the frame transform induced by a closed loop:
\[
\Phi_{\gamma} = T_{\gamma}.
\]
A loop exhibits non-trivial holonomy when:
\[
\Phi_{\gamma} \neq I_n.
\]

\section{Fun-Sphere Trigger Condition}

A point $x \in M$ triggers Fun-Sphere formation under loop $\gamma$ when:
\begin{enumerate}
    \item The loop returns to the same point:
    \[
    P_{\gamma}(x) = x.
    \]
    \item The frame transforms:
    \[
    \Phi_{\gamma} \neq I_n.
    \]
\end{enumerate}

\section{Fun-Sphere Definition}

The Fun-Sphere is the holonomy-invariant region:
\[
\mathcal{F} =
\left\{
x \in M \mid
P_{\gamma}(x) = x \ \text{and} \ \Phi_{\gamma} \neq I_n
\right\}.
\]

\section{Expansion Procedure}

Starting from $x_0$, the Fun-Sphere expands via:
\[
x_{k+1} = P_{\gamma}(x_k).
\]
Points are added whenever the trigger condition holds.

\section{Narrative Ecology}

Narrative evolution is modeled as reversible operator flow:
\[
x_{t+1} = T x_t,
\qquad
x_{t-1} = T^{-1} x_t.
\]
A reversible narrative loop is:
\[
x_{\text{forward}} = T_{\gamma} x_0,
\qquad
x_{\text{backward}} = T_{\gamma}^{-1} x_{\text{forward}}.
\]

\section{Interpretation}

This toy model demonstrates:
\begin{itemize}
    \item narrative as manifold traversal,
    \item story actions as operators,
    \item plot twists as holonomy,
    \item reversible story physics via operator inverses,
    \item pocket universes as holonomy-invariant submanifolds.
\end{itemize}

\section{Future Work}

Future extensions include:
\begin{itemize}
    \item nonlinear operator families,
    \item curvature-aware manifolds,
    \item multi-sphere holonomy networks,
    \item spectral narrative operators,
    \item geometric AI interpretability via holonomy signatures.
\end{itemize}

\bibliographystyle{plain}
\bibliography{preprint_bibliography}

\end{document}
